\section{Actual deep roots and multiplicative phase coherence}
\label{mc-section}

Let \(n\geq7\) be prime, \(B=n^4\), and let \(k\) range over the
integers \(B\leq k<2B\). Put
\[
 a_k=3k,\quad w_k=a_k+\zeta,\quad Q_k=a_k^2+a_k+1,\qquad
 T_n(k)=|P(w_k^n)|,\quad t_k=\log c_n(k),\quad Y=6n^3.
\]
The elements and their powers are primitive and unramified.
Here \(P(A+B\zeta)=AB(A+B)\), and \(c_n(k)\) is the largest
side after positive-sector normalization. Thus \(n<Y<B\).
In the pair arguments, \(a,b,c,d\) denote root parameters \(a_k\),
not the three entries of the resulting ABC triple.

\begin{lemma}[A finite group at actual fourth depth]\label{mc-torsion}
For a prime \(q>Y\), put
\[
 \mathcal{A}_q=\{k:B\leq k<2B,\ q^4\mid T_n(k)\}.
\]
For each such root, \(Q_k\) is a unit modulo \(q\), and
\[
 \alpha_k=\frac{a_k+\zeta}{a_k+\bar\zeta}
 \quad\hbox{belongs to}\quad
 H_q=\{x\in((\mathbb Z/q^4\mathbb Z)[\zeta])^\times:
                 x\bar x=1,\ x^{3n}=1\}.
\]
One has \(|H_q|\leq3n\).
\end{lemma}
\begin{proof}
The exact identity
\[
 z^3-\bar z^3=3\sqrt{-3}\,P(z)
\]
has a unit scalar at \(q>3\). If \(q\mid Q_k\), in the inert case
the reduction of \(w_k\) would be zero, contrary to its second
coordinate being one. In the split case exactly one component is
zero and the other is not, so the cubes of \(w_k^n\) and
\(\bar w_k^n\) could not coincide. Thus a boundary hit implies
\(q\nmid Q_k\), and the identity proves the stated group membership.

In the split case the quadratic ring is a product of two copies of
\(\mathbb Z/q^e\mathbb Z\), and a norm-one element is \((u,u^{-1})\).
Modulo \(q\), at most \(3n\) elements are killed by \(3n\).
Each lifts uniquely at every depth because the derivative of
\(X^{3n}-1\) is a unit at a root. In the inert case the residue
norm-one group is cyclic of order \(q+1\). It too has at most
\(3n\) such elements. Substituting \(x+q^e y\) gives a unique
linear lifting equation over \(\mathbb F_{q^2}\), with invertible
derivative. Reduction therefore injects \(H_q\) into a set of size
at most \(3n\). This counts a finite group; no distribution
assumption is made for the actual roots.
\end{proof}

\begin{lemma}[Integer phase rigidity and positive divisor fibers]
\label{mc-phase}
Two ordered pairs of roots in \(\mathcal{A}_q\) having the same product
in \(H_q\) have the same actual rational phase:
\[
 \frac{ab-1}{a+b+1}=\frac{cd-1}{c+d+1}.
\]
For a fixed nonempty phase \(u/v\), in lowest positive terms, its
ordered-pair fiber in the whole block has at most
\(\tau(u^2+uv+v^2)\) members, and
\[
 1\leq u^2+uv+v^2\leq2500B^4.
\]
\end{lemma}
\begin{proof}
The actual product coordinates are
\[
 (a+\zeta)(b+\zeta)=(ab-1)+(a+b+1)\zeta.
\]
Their positive coordinates satisfy \(ab-1\leq36B^2\) and
\(a+b+1\leq12B\). Equality of the two ratio products modulo
\(q^4\), after multiplication by unit denominators, implies
\[
 q^4\mid D,\qquad
 D=(ab-1)(c+d+1)-(cd-1)(a+b+1).
\]
The conjugate cross difference is a unit multiple of this integer.
But
\[
 |D|\leq864B^3<1296B^3<q^4
\]
because \(q>6B^{3/4}\), so \(D=0\).

Reducing the positive phase gives \(u\leq36B^2\), \(v\leq12B\).
For any pair with that phase,
\begin{equation}\label{mc-divisor-fiber}
 (va-u)(vb-u)=u^2+uv+v^2=:K.
\end{equation}
Both factors are positive, since
\((ab-1)/(a+b+1)<\min(a,b)\).
The first factor is a positive divisor of \(K\), determines \(a\),
and then determines \(b\). This is an injection of the actual
ordered fiber into the positive divisors, including diagonal and
exchanged pairs. Finally
\[
 K\leq1296B^4+432B^3+144B^2\leq2500B^4.
\]
\end{proof}

\begin{theorem}[A uniform count of actual deep roots]\label{mc-count}
For every \(\varepsilon>0\) there is a constant \(C_\varepsilon\),
independent of \(n,q\), such that
\[
 |\mathcal{A}_q|\leq C_\varepsilon\sqrt n\,B^{2\varepsilon}.
\]
If \(\mathcal{A}_q\ne\varnothing\), its prime \(q\) has exact
homogeneous rank \(n\) and \(q\equiv1\) or \(-1\pmod{3n}\).
\end{theorem}
\begin{proof}
For every \(m\geq1\), the elementary divisor estimate is
\(\tau(m)\leq C_\varepsilon m^\varepsilon\).
Indeed for \(p\geq2^{1/\varepsilon}\) and \(j\geq0\),
\(j+1\leq2^j\leq p^{\varepsilon j}\).
For each of the finitely many smaller primes,
\(\sup_{j\geq0}(j+1)p^{-\varepsilon j}\) is finite. Multiplication
over prime powers proves the estimate.

All \(|\mathcal{A}_q|^2\) ordered pairs map into \(H_q\).
Each occupied image has one actual phase by Lemma~\ref{mc-phase},
and at most \(C_\varepsilon(2500B^4)^\varepsilon\) pairs.
There are at most \(3n\) images. Taking square roots proves the
bound, including the empty-set case.

The rank of \(((w_k/\bar w_k)^3\bmod q)\) divides the prime \(n\).
Rank one would give
\(v_q(T_n(k))=v_q(T_1(k))\), since \(q>Y>n\).
Here \(T_1(k)=a_k(a_k+1)\) has two coprime factors, each below
\(6B+1\), whereas \(q^2>36n^6>6B+1\). Rank one cannot even
have depth two. The rank is therefore \(n\).
The norm-one group gives \(n\mid q-\chi_q\), where
\(\chi_q=1\) in the split case and \(-1\) in the inert case.
Also \(3\mid q-\chi_q\), and \(\gcd(3,n)=1\); hence the stated
two progressions.
\end{proof}

\begin{theorem}[A simultaneous interval of depth at most three]
\label{mc-window}
For real \(Z>Y\), let \(\mathcal{E}(n,Z)\) be the actual set of indices
with some prime \(Y<q\leq Z\) of boundary depth at least four. Then
\[
 \frac{|\mathcal{E}(n,Z)|}{B}
 \leq
 \frac{C_\varepsilon ZB^{2\varepsilon}}
      {B\sqrt n\,\log(Z/(3n))}.
\]
Fix \(0<\delta<1/2\). For all sufficiently large prime \(n\), put
\(Z_\delta=\lfloor Bn^{1/2-\delta}\rfloor\). Then
\begin{equation}\label{mc-exception}
 \frac{|\mathcal{E}(n,Z_\delta)|}{B}
       =O_\delta\!\left(\frac{n^{-\delta/2}}{\log n}\right).
\end{equation}
Outside this one set, every supported prime in the entire interval
\((Y,Z_\delta]\) has depth at most three.
\end{theorem}
\begin{proof}
The two-progression Brun--Titchmarsh estimate used in
Section~\ref{sw-section} gives at most
\[
 \frac{4Z}{\varphi(3n)\log(Z/(3n))}
 =\frac{2Z}{(n-1)\log(Z/(3n))}
\]
eligible primes. Apply Theorem~\ref{mc-count} at every such prime
and take a union of actual index sets. This proves the first bound,
with \(n/(n-1)\) absorbed into the constant.
For the second take \(\varepsilon=\delta/16\).
Since \(B^{2\varepsilon}=n^{8\varepsilon}\), the remaining power is
\(-\delta+8\varepsilon=-\delta/2\); the denominator is asymptotic
to \((7/2-\delta)\log n\). The floor has no effect on this limit.
\end{proof}

Write
\begin{gather*}
 J(T)=\sum_{q\mid T}(v_q(T)-3)\log q,\qquad
 L_Y(T)=\sum_{\substack{q\leq Y\\q\mid T}}v_q(T)\log q,\\
 W_\delta(k)=\sum_{\substack{q>Z_\delta\\q\mid T_n(k)}}
                  (v_q(T_n(k))-3)\log q .
\end{gather*}
For every root outside \(\mathcal{E}(n,Z_\delta)\), the intermediate
signed cost is exactly
\[
 -2\!\sum_{\substack{Y<q\leq Z_\delta\\v_q(T_n)=1}}\log q
 -\!\sum_{\substack{Y<q\leq Z_\delta\\v_q(T_n)=2}}\log q\leq0.
\]
Thus \(J(T_n)\leq W_\delta+L_Y(T_n)\), without assuming that the
full mass in this interval is small.

Theorem~\ref{fml-mean} and the elementary fixed-prime estimates in
Section~\ref{ebi-section}, evaluated at \(Y=6n^3\), give
\[
 \frac1B\sum_k\frac{L_Y(T_n(k))}{t_k}
       =O\!\left(\frac{\log n}{n}\right).
\]
The harmonic term is \(O(\log n/n)\), the endpoint is
\(O(1/(n\log n))\), and lifting with primes two and three costs
\(O(1/n)\). Put \(\eta_n=\sqrt{\log n/n}\).
Markov's inequality and \eqref{mc-exception} give one actual set of
relative size at least
\[
 1-C_\delta\left(\eta_n+\frac{n^{-\delta/2}}{\log n}\right)
\]
on which
\begin{align}
 \frac{J(T_n)}{t_k}&\leq\frac{W_\delta}{t_k}+\eta_n,
       \label{mc-signed}\\
 \frac{\log\operatorname{rad}(T_n)}{t_k}
   &\geq1-\frac4{3n}-\frac{\eta_n}{3}
                    -\frac{(W_\delta)_+}{3t_k}.
       \label{mc-radical}
\end{align}
The second exceptional term is not absorbed into \(\eta_n\).
The last inequality uses
\(\log T_n/t_k\geq3-4/n\) from Lemma~\ref{ebi-angle}.
A sequence in these sets with \((W_\delta)_+/t_k\to0\) satisfies
the restricted ABC estimate at every fixed tolerance eventually.
The far condition has not been proved.

The new mechanism is that congruences at an actual fourth depth
become integer phase identities with finite divisor fibers.
The fifteen declarations of \texttt{DeepRootPhase\allowbreak Arithmetic.lean}
formalize the actual product identities, positivity, injection into
\texttt{Nat.divisors}, the fiber cardinality bound, the determinant
height bound, and large-modulus rigidity. Its independent source
reviews and fresh Lean~4.32.0 verification are recorded separately.
Finite-ring torsion, lifting, the asymptotic divisor estimate,
Brun--Titchmarsh, and the final analytic consequences remain ordinary
proofs. No assertion that all those steps are formalized is made.

The interval has nonpositive signed cost on the stated majority,
but its negative credit has no proved positive proportion of the
height. Its full mass has not been shown negligible at this larger
cutoff, so the earlier almost-full far-mass statement is not
transferred here. The farther private top-rank primes and the
exceptional roots remain unresolved.
